%%%%%%%%%%  scan idx 89  |  folio 78  %%%%%%%%%%

\underline{Proposición 3.7.}

\begin{enumerate}[label=(\alph*)]
\item Sean $A$ y $S$ $R$-álgebras. Si
\[
 \operatorname{Tor}_i^R(A,S)=0\qquad(0<i\leq n),
\]
entonces, para todo haz $F$ sobre
$\underline C_{A\otimes_RS|S}$, la aplicación canónica
\[
 H^i(A\otimes_RS|S,F)
 \longrightarrow H^i(A|R,\widetilde SF)
\]
es un isomorfismo para todo $i\leq n$. Aquí
\[
 S:\underline C_{A|R}\longrightarrow
 \underline C_{A\otimes_RS|S},
 \qquad X\longmapsto X\otimes_RS.
\]

\item Sea $P$ un álgebra polinómica sobre $R$. Entonces, para todo haz
aditivo $F$ sobre $\underline C_{A\otimes_RP|R}$, la aplicación
canónica
\[
 H^n(A\otimes_RP|R,F)\longrightarrow H^n(A|R,F)
\]
es un isomorfismo para todo $n>0$.

\item Sean $S\longrightarrow R\longrightarrow A$ homomorfismos de
anillos. Entonces, para todo haz aditivo $F$ sobre
$\underline C_{A|S}$, hay una sucesión exacta
\begin{multline*}
0\longrightarrow H^0(A|R,F_R)\longrightarrow H^0(A|S,F)
\longrightarrow H^0(R|S,F)\longrightarrow H^1(A|R,F_R)\\
\longrightarrow H^1(A|S,F)\longrightarrow H^1(R|S,F)
\longrightarrow\cdots\longrightarrow H^n(A|R,F_R)\\
\longrightarrow H^n(A|S,F)\longrightarrow H^n(R|S,F)
\longrightarrow H^{n+1}(A|R,F_R)\longrightarrow\cdots,
\end{multline*}
donde $F_R$ es el haz sobre $\underline C_{A|R}$ definido por
\[
 F_R(X)=\ker\bigl(F(X)\longrightarrow F(R)\bigr)
\]
para todo $X\in\operatorname{Ob}(\underline C_{A|R})$, mediante el
funtor de olvido
$\underline C_{A|R}\longrightarrow\underline C_{A|S}$.
\end{enumerate}

\underline{Demostración.}

(a) Considérese el funtor continuo
\[
 S:\underline C_{A|R}\longrightarrow
 \underline C_{A\otimes_RS|S},
 \qquad S(X)=X\otimes_RS.
\]
Como $S$ envía las álgebras polinómicas sobre $R$ a álgebras
polinómicas sobre $S$, el funtor
\[
 \widetilde S:\widetilde{\underline C}_{A\otimes_RS|S}
 \longrightarrow\widetilde{\underline C}_{A|R}
\]
es exacto por 3.3(c). Se obtiene, por tanto, la aplicación canónica
\[
 \check S\,\underline H^{\bullet}_{A\otimes_RS|S}(F)
 \longrightarrow
 \underline H^{\bullet}_{A|R}(\widetilde SF)
\]
y el siguiente diagrama conmutativo de sucesiones espectrales:
\begin{center}
\small
\begin{tikzcd}[column sep=huge,row sep=large]
E_S^{p,q}(S(X))=
 \check H^p\bigl(X\otimes_RS|S,
 \underline H^q_{A\otimes_RS|S}(F)\bigr)
 \arrow[r,Rightarrow] \arrow[d]
& H^{p+q}(X\otimes_RS|S,F) \arrow[dd]\\
\check H^p\bigl(X|R,
 \check S\,\underline H^q_{A\otimes_RS|S}(F)\bigr)
 \arrow[d] & \\
E_R^{p,q}(X)=
 \check H^p\bigl(X|R,
 \underline H^q_{A|R}(\widetilde SF)\bigr)
 \arrow[r,Rightarrow]
& H^{p+q}(X|R,\widetilde SF).
\end{tikzcd}
\end{center}

%%%%%%%%%%  scan idx 90  |  folio 79  %%%%%%%%%%

Por 3.6(a), la aplicación
\[
 E_S^{p,q}(S(X))\longrightarrow
 \check H^p\bigl(X|R,
 \check S\,\underline H^q_{A\otimes_RS|S}(F)\bigr)
\]
es un isomorfismo siempre que
$\operatorname{Tor}_1^R(X,S)=0$. Supongamos inductivamente que
\[
 [\check S\,\underline H^q_{A\otimes_RS|S}(F)](X)
 \longrightarrow
 [\underline H^q_{A|R}(\widetilde SF)](X)
\]
es un isomorfismo para todo $q<n$ y todo
$X\in\operatorname{Ob}(\underline C_{A|R})$ que satisface
\[
 \operatorname{Tor}_i^R(X,S)=0\qquad(0<i\sisi{\leq}{<} n).
\]
Se sigue de 3.6(b) que
$E_S^{p,q}(S(X))\longrightarrow E_R^{p,q}(X)$ es un isomorfismo
para todo $q<n$ y todo $p$, bajo la misma hipótesis de Tor. Además,
$E_S^{0,n}(S(X))\longrightarrow E_R^{0,n}(X)$ es un isomorfismo,
puesto que ambos miembros son cero. Por tanto,
\[
 H^n(X\otimes_RS|S,F)
 \longrightarrow H^n(X|R,\widetilde SF)
\]
es un isomorfismo siempre que
$\operatorname{Tor}_i^R(X,S)=0$ para $0<i\leq n$. Equivalentemente,
\[
 [\check S\,\underline H^i_{A\otimes_RS|S}(F)](X)
 \longrightarrow
 [\underline H^i_{A|R}(\widetilde SF)](X)
\]
es un isomorfismo para todo $i\leq n$ bajo esa hipótesis.

(b) Sea
\[
 f:\underline C_{A|R}\longrightarrow
 \underline C_{A\otimes_RP|R},
 \qquad f(X)=X\otimes_RP.
\]
Elíjase una flecha suprayectiva $P'\longrightarrow X$ en
$\underline C_{A|R}$, siendo $P'$ un álgebra polinómica sobre $R$.
Entonces
\[
 P'\otimes_RP\longrightarrow X\otimes_RP
\]
es una flecha suprayectiva en $\underline C_{A\otimes_RP|R}$, y
$P'\otimes_RP$ es de nuevo un álgebra polinómica sobre $R$. Puesto que
$P$ es plano
sobre $R$, 3.5(i) da
\begin{align*}
 \check H^{\bullet}(X|R,\check fG)
 &=H^{\bullet}(\{P'\longrightarrow X\},\check fG)\\
 &=H^{\bullet}(\{P'\otimes_RP\longrightarrow X\otimes_RP\},G)\\
 &=\check H^{\bullet}(X\otimes_RP|R,G)
\end{align*}
para todo
$G\in\operatorname{Ob}(\overset{\vee}{\underline C}_{A\otimes_RP|R})$.
Así, $\widetilde f$ envía los haces acíclicos a haces acíclicos. Por otra
parte, $f$ envía las álgebras polinómicas sobre $R$ a álgebras
polinómicas sobre $R$, de modo que
\[
 \widetilde f:\widetilde{\underline C}_{A\otimes_RP|R}
 \longrightarrow\widetilde{\underline C}_{A|R}
\]
es exacto por 3.3(c). Se sigue que
\[
 H^{\bullet}(A\otimes_RP|R,F)
 \longrightarrow H^{\bullet}(A|R,\widetilde fF)
\]
es un isomorfismo para todo
$F\in\operatorname{Ob}(\widetilde{\underline C}_{A\otimes_RP|R})$.

Para cada $X\in\operatorname{Ob}(\underline C_{A|R})$ hay una aplicación
canónica
\[
 (\widetilde fF)(X)=F(X\otimes_RP)
 \longrightarrow F(X)\oplus F(P),
\]
y por tanto una aplicación canónica
$\widetilde fF\longrightarrow F\oplus F(P)$, donde $F(P)$ denota el haz
constante y el $F$ del miembro derecho denota la restricción de $F$ a
$\underline C_{A|R}$. Si $F$ es aditivo, esta aplicación es un
isomorfismo. Por consiguiente, para todo $n>0$,
\[
 H^n(A|R,\widetilde fF)
 \xrightarrow{\ \sim\ }
 H^n(A|R,F\oplus F(P))=H^n(A|R,F),
\]
porque el haz constante $F(P)$ es flácido. Esto demuestra (b).

%%%%%%%%%%  scan idx 91  |  folio 80  %%%%%%%%%%

(c) Este argumento fue comunicado por Rinehart. Considérese el funtor de
olvido
$\beta:\underline C_{A|R}\longrightarrow\underline C_{A|S}$.
Hay una sucesión espectral
\[
 E^{p,q}=H^p\bigl(A|R,
 \beta^{\#}\underline H^q_{A|S}(F)\bigr)
 \Longrightarrow H^{p+q}(A|S,F),
\]
donde
\[
 \beta^{\#}=G_{A|R}\circ\beta^{\vee},
 \qquad
 G_{A|R}:\overset{\vee}{\underline C}_{A|R}
 \longrightarrow\widetilde{\underline C}_{A|R}
\]
es el funtor de haz asociado.

Supongamos ahora que $F$ es aditivo. Afirmamos que
$\beta^{\#}\underline H^q_{A|S}(F)$ es un haz constante para todo
$q>0$. En efecto, si $P$ es un álgebra polinómica sobre $R$, entonces
$P=R\otimes_SP'$ para un álgebra polinómica $P'$ sobre $S$, y
\[
 [\beta^{\#}\underline H^q_{A|S}(F)](P)
 =H^q(P|S,F)
 =H^q(R\otimes_SP'|S,F)
 \longrightarrow H^q(R|S,F)
\]
es un isomorfismo para $q>0$ por (b), y es suprayectiva para $q=0$.
Por 3.3(c), la aplicación canónica
\[
 \beta^{\#}\underline H^q_{A|S}(F)
 \longrightarrow H^q(R|S,F)
\]
al haz constante correspondiente es, por tanto, un isomorfismo para
$q>0$ y una suprayección para $q=0$. De aquí que $E^{p,q}=0$ siempre que
$p,q>0$, y la sucesión espectral proporciona la sucesión exacta
\[
 (*)\quad
 0\longrightarrow E^{1,0}\longrightarrow H^1
 \longrightarrow E^{0,1}\longrightarrow E^{2,0}
 \longrightarrow H^2\longrightarrow E^{0,2}
 \longrightarrow E^{3,0}\longrightarrow H^3\longrightarrow\cdots.
\]
Aquí
\[
 E^{0,q}=[\beta^{\#}\underline H^q_{A|S}(F)](A)
 \xrightarrow{\ \sim\ }H^q(R|S,F)\quad(q>0),
 \qquad
 E^{p,0}=H^p(A|R,\widetilde\beta F).
\]
Por otra parte, la sucesión exacta
\[
 0\longrightarrow F_R\longrightarrow\widetilde\beta F
 \longrightarrow F(R)\longrightarrow0
\]
da el diagrama exacto
\begin{center}
\scriptsize
\begin{tikzcd}[column sep=small,row sep=large]
0 \arrow[r]
& H^0(A|R,F_R) \arrow[r]
& H^0(A|R,\widetilde\beta F) \arrow[r] \arrow[d,equal]
& H^0(A|R,F(R)) \arrow[r] \arrow[d,equal]
& H^1(A|R,F_R) \arrow[r]
& H^1(A|R,\widetilde\beta F) \arrow[r]
& 0\\
&& H^0(A|S,F) & H^0(R|S,F) &&&
\end{tikzcd}
\end{center}
etiquetado $(**)$ en la fuente. La combinación de $(*)$ y $(**)$ da la
sucesión exacta buscada. \hfill$\square$
